\section{Conclusion}

We investigated whether open-source chest X-ray VLMs can be selectively trusted under distribution shift, and whether their apparent reliability is grounded in visual evidence. Across seven models spanning six families, visual reliance grades continuously from image-dependent to image-blind rather than splitting into a binary, and the image-blind models compound low visual grounding with a default-normal prediction strategy while still emitting structurally valid outputs and apparently favourable selective metrics. Under image-quality shift, degradation does not erode macro accuracy; it selectively suppresses the positive recall of the minority pathology labels, and this effect is degradation-type-specific, with reduced resolution and added noise the dominant factors. The models additionally separate into three safety-failure modes -- default-normal, refusal-to-answer, and prose-without-the-contract -- none of which requires a parse error.

These findings show that accuracy, calibration, or selective risk alone mischaracterise CXR VLM trustworthiness, and that reliability must be assessed jointly through visual grounding, positive recall, and failure mode. The shift-aware selective trust framework described here -- which conditions accept/abstain decisions on these signals -- is offered as a conceptual blueprint grounded in the empirical results, whose predictive components await ablation and whose pre-run observations await the full multi-centre run. We close by flagging the two hypotheses that run is designed to arbitrate: that the degradation-type ranking holds at scale, and that conditioning selective trust on visual reliance and positive recall improves risk--coverage over confidence alone.